% GENERATED FILE. Edit canonical lessons, then run npm run generate:books.
% canonical-source-hash: fnv1a64:ad0cf9f716301792
% canonical-lessons: LA-C286-navis, LA-C286-naulum, LA-C286-amitto, LA-R286-by-sea

\chapter{By Sea}
\label{ch:la-by-sea}
\hlchaptermodality{286}

\begin{chapteropening}
\textbf{By the end of this chapter:} \emph{I can book a passage on a ship, pay the fare, and report that something went wrong on the way.}

Complete LA-R286-by-sea: I can book a passage on a ship, pay the fare, and report that something went wrong on the way.
\end{chapteropening}

\section[nāvis, nāvis]{nāvis, nāvis — a ship}

\label{lesson:LA-C286-navis}



\begin{quote}
Say \textquotedblleft{}I am making a journey\textquotedblright{}. (\emph{Iter faciō.}) Then \textquotedblleft{}to sail\textquotedblright{}, and \textquotedblleft{}a sailor\textquotedblright{}. (\emph{Nāvigāre. Nauta.})

A long Roman journey often went by sea. The word for what the sailor sails has been hiding inside \emph{nāvigāre} all along.
\end{quote}

\subsection*{You'll want to know: nāvis, nāvis}
\textbf{nāvis, nāvis} — \textquotedblleft{}a ship\textquotedblright{}. It is feminine, and its accusative is \textbf{nāvem}.

\begin{quote}
\textbf{In nāvem ascendō.} — I go on board the ship. \textbf{Nāvis solvit.} — The ship sets sail.
\end{quote}

You own \textbf{ascendō}, \textquotedblleft{}to climb, to go up\textquotedblright{}, and Romans went \emph{up} into a ship. And you own \textbf{solvō}, \textquotedblleft{}to untie\textquotedblright{}: a ship that \emph{solvit} has cast off its ropes from the shore.

\begin{cousinweb}
\textbf{nāvis} is inherited, usually traced to Proto-Indo-European \emph{*néh\textsubscript{2}us}, \textquotedblleft{}boat\textquotedblright{} (EtymDB 2.1). Greek \emph{naus}, \textquotedblleft{}ship\textquotedblright{}, is a cousin (Lewis \& Short set them side by side). From the Latin come English \emph{naval}, \emph{navy} and the \emph{nave} of a church, and \emph{nāvigō} is \emph{nāvis} with \emph{agō}, \textquotedblleft{}to drive\textquotedblright{} (Lewis \& Short): to drive a ship.
\end{cousinweb}

\subsection*{Your turn}
Say these aloud:

\begin{itemize}
\raggedright
  \item \emph{nāvis, nāvis}
  \item \emph{In nāvem ascendō.}
  \item \emph{Nāvis solvit.}
  \item \emph{Herī nāvis solvit, et diū nāvigāvimus.}
  \item the journey and the ship — \emph{iter}, then \emph{nāvis}
\end{itemize}

\begin{tcolorbox}[breakable,colback=teal!4,colframe=teal!35!black,title={Before you move on}]
What does \emph{nāvis, nāvis} mean? (\textbf{A ship.}) Say \textquotedblleft{}I go on board\textquotedblright{}. (\emph{\textbf{In nāvem ascendō.}}) Say \textquotedblleft{}the ship sets sail\textquotedblright{}. (\emph{\textbf{Nāvis solvit.}}) Which Greek word is its cousin? (\emph{\textbf{Naus.}})
\end{tcolorbox}
\section[naulum, naulī]{naulum, naulī — the fare}

\label{lesson:LA-C286-naulum}



\begin{quote}
Say \textquotedblleft{}I go on board the ship\textquotedblright{}. (\emph{In nāvem ascendō.})

Romans had no tickets. A passenger found a ship going his way and paid its master a fare.
\end{quote}

\subsection*{You'll want to know: naulum, naulī}
\textbf{naulum, naulī} — \textquotedblleft{}a fare, passage money\textquotedblright{}. It is neuter, like \emph{pretium}, \textquotedblleft{}a price\textquotedblright{}, which is the word for what it costs.

\begin{quote}
\textbf{Naulum solvō.} — I pay the fare. \textbf{Decem dēnāriōs solvī.} — I paid ten denarii.
\end{quote}

\emph{Solvō} does both jobs of this chapter: it unties a ship from the shore, and it pays what you owe.

The man you pay is the \textbf{magister nāvis}, \textquotedblleft{}the master of the ship\textquotedblright{}: the same \emph{magister} as a teacher, the one in charge.

\begin{cousinweb}
\textbf{naulum} was borrowed from Greek \emph{naulon}, \textquotedblleft{}passage money\textquotedblright{} (Lewis \& Short), a word usually connected with \emph{naus}, \textquotedblleft{}ship\textquotedblright{}, the Greek cousin of \emph{nāvis}. The satirist Juvenal uses it, and so do the Roman law books, for the money owed to a ship; the law books also name the \emph{magister nāvis}, the ship's master (Lewis \& Short).
\end{cousinweb}

\subsection*{Your turn}
Say these aloud:

\begin{itemize}
\raggedright
  \item \emph{naulum, naulī}
  \item \emph{Naulum solvō.}
  \item \emph{Prīmum naulum solvī; deinde in nāvem ascendī.}
  \item \emph{Magister nāvis decem dēnāriōs accēpit.} — The ship's master took ten denarii
  \item on land — \emph{raeda}, a travelling carriage — then by sea: \emph{nāvis} and \emph{naulum}
\end{itemize}

\begin{tcolorbox}[breakable,colback=teal!4,colframe=teal!35!black,title={Before you move on}]
What does \emph{naulum, naulī} mean? (\textbf{A fare}, for a ship.) Say \textquotedblleft{}I pay the fare\textquotedblright{}. (\emph{\textbf{Naulum solvō.}}) Who takes it? (The \emph{\textbf{magister nāvis}}.) Which language lent Latin the word? (\textbf{Greek.})
\end{tcolorbox}
\section[āmittō, āmittere]{āmittō — something went wrong}

\label{lesson:LA-C286-amitto}



\begin{quote}
Ask \textquotedblleft{}what happened?\textquotedblright{} (\emph{Quid accidit?}) Then say \textquotedblleft{}I paid the fare\textquotedblright{}. (\emph{Naulum solvī.})

On a journey, the answer to \emph{quid accidit} is often bad news.
\end{quote}

\subsection*{You'll want to know: āmittō, āmittere}
\textbf{āmittō, āmittere} — \textquotedblleft{}to lose, to let slip\textquotedblright{}. Its perfect is \textbf{āmīsī}, \textquotedblleft{}I have lost\textquotedblright{}.

\begin{quote}
\textbf{Sarcinam āmīsī!} — I have lost my pack! \textbf{Occāsiōnem āmīsī.} — I let the chance slip.
\end{quote}

You own \textbf{perdō}, \textquotedblleft{}to lose\textquotedblright{}, too, and the two are not the same. \emph{Perdō} is losing by wasting or wrecking: money thrown away, a thing ruined. \emph{Āmittō} is losing by letting go: it slipped away from you.

\begin{grammarlens}[title={Grammar Lens: reporting the problem}]
A problem on the road is an event, so it comes in the \textbf{perfect}:

\begin{quote}
\textbf{— Quid accidit?} — What happened? \textbf{— Nāvis iam discessit!} — The ship has already left!
\end{quote}

\emph{Discessit} is the perfect of \emph{discēdō}, \textquotedblleft{}to go away\textquotedblright{}, and \emph{iam}, \textquotedblleft{}already\textquotedblright{}, makes it the missed boat. Answer \emph{quid accidit} with what happened, then say what it means for you: \emph{Sarcinam āmīsī; nummōs nōn habeō.} I have lost my pack; I have no money.
\end{grammarlens}

\begin{cousinweb}
\textbf{āmittō} is built from \emph{ab}, \textquotedblleft{}away\textquotedblright{}, and \emph{mittō}, \textquotedblleft{}to send\textquotedblright{} (EtymDB 2.1; Lewis \& Short): to send away from yourself. In early Latin it meant \textquotedblleft{}to let someone go\textquotedblright{}, and from there \textquotedblleft{}to let a thing slip\textquotedblright{}, and so \textquotedblleft{}to lose\textquotedblright{}. Its perfect \emph{āmīsī} keeps the perfect of \emph{mittō}, \emph{mīsī}.
\end{cousinweb}

\subsection*{Your turn}
Say these aloud:

\begin{itemize}
\raggedright
  \item \emph{āmittō, āmittere} — then \emph{āmīsī}
  \item \emph{Sarcinam āmīsī!}
  \item \emph{Quid accidit? — Nāvis iam discessit!}
  \item \emph{Naulum solvī, sed nāvis iam discessit.}
  \item a problem of your own on a journey, and what it means for you
\end{itemize}

\begin{tcolorbox}[breakable,colback=teal!4,colframe=teal!35!black,title={Before you move on}]
What does \emph{āmittō, āmittere} mean? (\textbf{To lose, to let slip.}) Say \textquotedblleft{}I have lost my pack\textquotedblright{}. (\emph{\textbf{Sarcinam āmīsī.}}) Say \textquotedblleft{}the ship has already left\textquotedblright{}. (\emph{\textbf{Nāvis iam discessit.}}) Which verb loses by wasting? (\emph{\textbf{Perdō.}})
\end{tcolorbox}
\section[(review)]{Review — a journey by sea}

\label{lesson:LA-R286-by-sea}



\begin{quote}
Say \textquotedblleft{}a ship\textquotedblright{}, then \textquotedblleft{}the fare\textquotedblright{}. (\emph{Nāvis. Naulum.})
\end{quote}

\begin{grammarlens}[title={What you've built}]
Two chapters, and a whole journey. The way through the town, \emph{quā viā …}, \emph{rēctā viā}, \emph{ad dextram}. The ship, the fare and the master who takes it. And the trouble: \emph{quid accidit?}, answered in the perfect.

\begin{quote}
\emph{Iter Rōmam faciēbam. Prīmum naulum solvī. Deinde sarcinam āmīsī. Postrēmō nāvis discessit!} — I was making a journey to Rome. First I paid the fare. Then I lost my pack. And last of all, the ship left!
\end{quote}

One line of scene, \emph{faciēbam}, and three events.
\end{grammarlens}

\subsection*{Your turn}
Say these aloud:

\begin{itemize}
\raggedright
  \item \emph{Quaesō, dīc mihi: quā viā ad nāvem?}
  \item the answer — \emph{Rēctā viā ī; deinde ad dextram.}
  \item \emph{Naulum solvō; deinde in nāvem ascendō.}
  \item \emph{Quid accidit? — Sarcinam āmīsī!}
  \item the whole journey in four lines: a scene, then three events
\end{itemize}

\begin{tcolorbox}[breakable,colback=teal!4,colframe=teal!35!black,title={Before you move on}]
Say \textquotedblleft{}I pay the fare\textquotedblright{}. (\emph{\textbf{Naulum solvō.}}) Say \textquotedblleft{}I have lost my pack\textquotedblright{}. (\emph{\textbf{Sarcinam āmīsī.}}) Ask \textquotedblleft{}what happened?\textquotedblright{} (\emph{\textbf{Quid accidit?}}) Say \textquotedblleft{}I am making a journey\textquotedblright{}. (\emph{\textbf{Iter faciō.}})
\end{tcolorbox}
